\documentclass[10pt, letterpaper]{article}

% Packages:
\usepackage[
    ignoreheadfoot, % set margins without considering header and footer
    top=2 cm, % seperation between body and page edge from the top
    bottom=2 cm, % seperation between body and page edge from the bottom
    left=2 cm, % seperation between body and page edge from the left
    right=2 cm, % seperation between body and page edge from the right
    footskip=1.0 cm, % seperation between body and footer
    % showframe % for debugging 
]{geometry} % for adjusting page geometry
\usepackage[explicit]{titlesec} % for customizing section titles
\usepackage{tabularx} % for making tables with fixed width columns
\usepackage{array} % tabularx requires this
\usepackage[dvipsnames]{xcolor} % for coloring text
\definecolor{primaryColor}{RGB}{0, 79, 144} % define primary color
\usepackage{enumitem} % for customizing lists
\usepackage{fontawesome5} % for using icons
\usepackage{amsmath} % for math
\usepackage[
    pdftitle={Chunxu Han's CV},
    pdfauthor={Chunxu Han},
    pdfcreator={LaTeX with RenderCV},
    colorlinks=true,
    urlcolor=primaryColor
]{hyperref} % for links, metadata and bookmarks
\usepackage[pscoord]{eso-pic} % for floating text on the page
\usepackage{calc} % for calculating lengths
\usepackage{bookmark} % for bookmarks
\usepackage{lastpage} % for getting the total number of pages
\usepackage{changepage} % for one column entries (adjustwidth environment)
\usepackage{paracol} % for two and three column entries
\usepackage{ifthen} % for conditional statements
\usepackage{needspace} % for avoiding page brake right after the section title
\usepackage{iftex} % check if engine is pdflatex, xetex or luatex
\usepackage{multicol}
\usepackage{datetime2}

% Ensure that generate pdf is machine readable/ATS parsable:
\ifPDFTeX
    \input{glyphtounicode}
    \pdfgentounicode=1
    \usepackage[T1]{fontenc}
    \usepackage[utf8]{inputenc}
    \usepackage{lmodern}
\fi

\usepackage[default, type1]{sourcesanspro} 

% Some settings:
\AtBeginEnvironment{adjustwidth}{\partopsep0pt} % remove space before adjustwidth environment
\pagestyle{empty} % no header or footer
\setcounter{secnumdepth}{0} % no section numbering
\setlength{\parindent}{0pt} % no indentation
\setlength{\topskip}{0pt} % no top skip
\setlength{\columnsep}{0.15cm} % set column seperation
\makeatletter
\let\ps@customFooterStyle\ps@plain % Copy the plain style to customFooterStyle
\patchcmd{\ps@customFooterStyle}{\thepage}{
    \color{gray}\textit{\small Chunxu Han - Page \thepage{} of \pageref*{LastPage}}
}{}{} % replace number by desired string
\makeatother
\pagestyle{customFooterStyle}

\titleformat{\section}{
    % avoid page braking right after the section title
    \needspace{4\baselineskip}
    % make the font size of the section title large and color it with the primary color
    \Large\color{primaryColor}
}{
}{
}{
    % print bold title, give 0.15 cm space and draw a line of 0.8 pt thickness
    % from the end of the title to the end of the body
    \textbf{#1}\hspace{0.15cm}\titlerule[0.8pt]\hspace{-0.1cm}
}[] % section title formatting

\titlespacing{\section}{
    % left space:
    -1pt
}{
    % top space:
    0.3 cm
}{
    % bottom space:
    0.2 cm
} % section title spacing

% \renewcommand\labelitemi{$\vcenter{\hbox{\small$\bullet$}}$} % custom bullet points
\newenvironment{highlights}{
    \begin{itemize}[
        topsep=0.10 cm,
        parsep=0.10 cm,
        partopsep=0pt,
        itemsep=0pt,
        leftmargin=0.4 cm + 10pt
    ]
}{
    \end{itemize}
} % new environment for highlights

\newenvironment{highlightsforbulletentries}{
    \begin{itemize}[
        topsep=0.10 cm,
        parsep=0.10 cm,
        partopsep=0pt,
        itemsep=0pt,
        leftmargin=10pt
    ]
}{
    \end{itemize}
} % new environment for highlights for bullet entries


\newenvironment{onecolentry}{
    \begin{adjustwidth}{
        0.2 cm + 0.00001 cm
    }{
        0.2 cm + 0.00001 cm
    }
}{
    \end{adjustwidth}
} % new environment for one column entries

\newenvironment{twocolentry}[2][]{
    \onecolentry
    \def\secondColumn{#2}
    \setcolumnwidth{\fill, 4.5 cm}
    \begin{paracol}{2}
}{
    \switchcolumn \raggedleft \secondColumn
    \end{paracol}
    \endonecolentry
} % new environment for two column entries

\newenvironment{threecolentry}[3][]{
    \onecolentry
    \def\thirdColumn{#3}
    \setcolumnwidth{1 cm, \fill, 4.5 cm}
    \begin{paracol}{3}
    {\raggedright #2} \switchcolumn
}{
    \switchcolumn \raggedleft \thirdColumn
    \end{paracol}
    \endonecolentry
} % new environment for three column entries

\newenvironment{header}{
    \setlength{\topsep}{0pt}\par\kern\topsep\centering\color{primaryColor}\linespread{1.5}
}{
    \par\kern\topsep
} % new environment for the header


% save the original href command in a new command:
\let\hrefWithoutArrow\href

% new command for external links:
\renewcommand{\href}[2]{\hrefWithoutArrow{#1}{\ifthenelse{\equal{#2}{}}{ }{#2 }\raisebox{.15ex}{\footnotesize \faExternalLink*}}}


\begin{document}
    \newcommand{\AND}{\unskip
        \cleaders\copy\ANDbox\hskip\wd\ANDbox
        \ignorespaces
    }
    \newsavebox\ANDbox
    \sbox\ANDbox{}

    \begin{header}
        \raggedright
        \fontsize{28 pt}{28 pt} 
        \textbf{Chunxu Han}

        \vspace{0.3 cm}

        \normalsize
        \mbox{{\footnotesize\faMapMarker*}\hspace*{0.13cm}Peder Skrams Gade 16A, København K, Denmark}%
        \kern 0.25 cm%
        \AND%
        \kern 0.25 cm%
        \mbox{\hrefWithoutArrow{mailto:s220311@dtu.dk}{{\footnotesize\faEnvelope[regular]}\hspace*{0.13cm}s220311@dtu.dk}}%
        \kern 0.25 cm%
        \AND%
        \kern 0.25 cm%
        \mbox{\hrefWithoutArrow{tel:+45 52 73 06 85}{{\footnotesize\faPhone*}\hspace*{0.13cm}+45 52 73 06 85}}%
        \kern 0.25 cm%
        \AND%
        \kern 0.25 cm%
        
        \mbox{\hrefWithoutArrow{https://www.linkedin.com/in/chunxu-han/}{{\footnotesize\faLinkedinIn}\hspace*{0.13cm}chunxu-han}}%
        \kern 0.25 cm%
        \AND%
        \kern 0.25 cm%
        \mbox{\hrefWithoutArrow{https://github.com/chunxubioinfor}{{\footnotesize\faGithub}\hspace*{0.13cm}chunxubioinfor}}%
    \end{header}

    \vspace{0.4 cm - 0.3 cm}


    % \section{Profile}



        
    %     \begin{onecolentry}
    %         I’m a Bioinformatics Master’s student with strong experience in statistical modelling and consulting.  
    %     \end{onecolentry}

    %     \vspace{0.2 cm}
        
    \section{Relevant Experience}

          
\begin{onecolentry}


\textbf{Quantitative Analysis \& Statistical Modelling} \\
I study bioinformatics, a field at the intersection of biology, statistics, and data science, and have two years of work experience at DTU HealthTech. I contributed to the development of the \textbf{R package} \href{https://github.com/kvittingseerup/IsoformSwitchAnalyzeR}{IsoformSwitchAnalyzeR}, including the introduction of a new statistical framework for isoform-switch analysis. Through this work, I gained hands-on experience in quantitative modelling, benchmarking, and data visualisation, which has contributed to a research paper currently under review.
        \vspace{0.3 cm}

\textbf{AI-Powered Pipeline Development} \\
I have also explored how state-of-the-art data science and machine learning methods can be applied to real-world problems. In my academic work, I have integrated machine learning and deep learning approaches into bioinformatics research workflows. Meanwhile, I have developed independent projects, including an AI-powered job search agent \href{https://github.com/chunxubioinfor/DailyJobMatch}{DailyJobMatch} for personalised job matching and a data pipeline for analysing trends in the Danish job market.

        \vspace{0.3 cm}
\textbf{Data Translation \& Consulting} \\
During my full-time internship at the marketing consulting firm Illuminera (IQVIA), I worked on \textbf{quantitative market research} projects, including \textbf{social media listening}, for pharmaceutical clients. I developed experience in transforming data and analysis into clear, actionable recommendations and in tailoring presentations and deliverables to different audiences and stakeholders. This strengthened both my consulting mindset and my ability to communicate analytical insights effectively.
\end{onecolentry}

    \section{Education}



        
        \begin{twocolentry}{
    Lyngby, Denmark

    Sept 2023 - Mar 2026
}
    \textbf{MSc in Bioinformatics and Systems Biology}, Technical University of Denmark
    \begin{highlights}
        \item Completed coursework in \textbf{Machine Learning} and \textbf{Deep Learning}, strengthening my ability to design data-driven analytical workflows and apply modern modelling approaches.
    \end{highlights}
\end{twocolentry}


% \begin{twocolentry}{
%     Beijing, China
    
%     Sept 2017 - June 2021
% }
%     \textbf{BSc in Biotechnology}, Beijing Normal University
%     \begin{highlights}
%         \item \textbf{Coursework:} Built a solid foundation in \textbf{Biotech} with courses covering ecology, molecular biology, cell biology, and bioinformatics, alongside laboratory training.
%     \end{highlights}
% \end{twocolentry}


    
    \section{Work Experience}

        \begin{twocolentry}{
    Målev, Denmark

    Sep 2025 – Mar 2026
}
    \textbf{Master Student}, Novo Nordisk

    \vspace{0.1cm}
    
     My master’s thesis focuses on developing computational frameworks to evaluate batch-to-batch variability in single-cell RNA-seq data:

    \begin{highlights}
        \item \textbf{Methodology development:} Developed Batch Mixing Index (BMI), an entropy-based metric that quantifies per-cell batch mixing and incorporates \textbf{permutation-based statistical testing}, as well as scBatch, a \textbf{Python package} integrating seven complementary metrics for manufacturing batch variability assessment.

        \vspace{0.1cm}

        \item \textbf{AI/ML approaches:} Explored state-of-the-art data science approaches by applying the Deep Learning model to embed cell therapy batches into a shared latent space for variability assessment.

        \vspace{0.1cm}

        % \item \textbf{Meetings and Collaboration:} Participate in team meetings to stay updated on current trends in the Pharma industry. Collaborate with cross-functional R\&D and CMC teams to ensure the biological interpretability and industrial relevance.
    \end{highlights}
\end{twocolentry}


        \vspace{0.3 cm}

        
        \begin{twocolentry}{
    Lyngby, Denmark

    Oct 2023 –  July 2025
}
    \textbf{Bioinformatics Research Assistant}, DTU HealthTech

    \vspace{0.1cm}
    
     I worked in Isoform Analysis Group for two years, which has been a transformative experience. I have grown from a student who only focuses on theoretical studies into a bioinformatician to deal with real-world problems:

    \begin{highlights}
        \item \textbf{Development and Maintenance of R Tools:} Contributed to developing and updating the \href{https://github.com/kvittingseerup/IsoformSwitchAnalyzeR}{IsoformSwitchAnalyzeR} R package. This experience sharpened my coding skills and can-do mindset. 

        \vspace{0.1cm}

        \item \textbf{Large-scale Data Analysis Pipelines:} Built and optimised processing pipelines for large-scale, high-dimensional data analyses, including automation, quality control, reproducible workflows, HPC-based parallel processing, and version control with Git.

        \vspace{0.1cm}

        \item \textbf{Cross-Team Analytical Support:} Collaborated daily with PhDs, postdocs, and master's, helping with data preprocessing and pipeline running. This strengthened my teamwork skills and equipped me with a \textbf{ collaborative mindset}. 

        % \vspace{0.1cm}

        % \item \textbf{Ad-hoc Work:} Managed various ad-hoc tasks, from literature reviews to data organization, learning how to \textbf{prioritize and organize} work effectively in a fast-paced research environment.
    \end{highlights}
\end{twocolentry}


        \vspace{0.3 cm}

        \begin{twocolentry}{
    Shanghai, China

    Feb 2023 – May 2023
}
    \textbf{Consulting Intern}, Illuminera, IQVIA

    \vspace{0.1cm}
    
    My first experience in a business setting, where I gained insights into the pharmaceutical market and developed professional skills:
    

    \begin{highlights}
        \item \textbf{Qualitative and Quantitative Research:} Conducted market trend and competitor analyses to develop actionable strategies for clients, while building a strong understanding of the pharmaceutical market in China and globally.

        \vspace{0.1cm}
 
        \item \textbf{Social Media Listening:} Led a social media listening project using platforms such as RedNote and online hospital communities, collecting and analysing large-scale user data to identify product gaps, map patient and consumer preferences, and generate strategic recommendations for \textbf{product positioning} and \textbf{targeted commercial campaigns}.

        \vspace{0.1cm}
        
        \item \textbf{Client and KOL Engagement:} Built professional relationships with clients and KOLs and coordinated with stakeholders to align project progress and delivery.

        \vspace{0.1cm}
        
        \item \textbf{Data-Driven Deliverables:} Utilized tools like PowerPoint and Power BI to deliver strategic reports, transforming complex data into actionable business insights. I learned how to formulate the business problems and storytelling to the clients. 
    \end{highlights}
\end{twocolentry}

        \vspace{0.3 cm}


    
%     \section{Publications}

        
%         \begin{onecolentry}
%             \begin{itemize}
%                 \item Jeppe Malthe Mikkelsen, \textbf{Chunxu Han}, Dimitrios S. Kanakoglou \& Kristoffer Vitting-Seerup. \textit{Isoform-level examination empowers analysis of both common and rare genetic variation}. \textbf{Genome Medicine}, Submitted.

%                 \vspace{0.1 cm}
    
%                 \item \textbf{Chunxu Han} \& Kristoffer Vitting-Seerup. \textit{IsoformSwitchAnalyzeR v2: Enabling Identification and Analysis of Isoform Switches with Functional Consequences from Long-read and Single-cell Sequencing Data}. In preparation.
% \end{itemize}
%         \end{onecolentry}

    \section{Hobby Projects}

          
\begin{onecolentry}


\textbf{Job Searching Agent} \\
Developed \href{https://github.com/chunxubioinfor/DailyJobMatch}{DailyJobMatch}, an automated job-matching agent that integrates web scraping, LLM-based scoring and workflow automation to collect, filter, and rank relevant job postings. Until now it has received 120+ stars on Github and around 5k likes on social media. I feel really rewarding to be able to leverage modern data science tools to real-world problems and contribute the community and help others.

        \vspace{0.3 cm}

\textbf{Bioinformatics Jobs Landscape analysis} \\
Developed a Danish bioinformatics job market analysis pipeline that scrapes and enriches 2,500+ job postings from jobindex.dk over a 5-year period. The project combines Playwright-based web scraping, GPT-4o for structured skill and sector extraction, and custom visualisations to map hiring trends, in-demand skills, and the rise of AI requirements across academia and industry. Built as a personal project to understand the job landscape before starting my own search.

\end{onecolentry}
    
    \section{Skills and Hobbies}

        \begin{onecolentry}
            \textbf{Programming:} Python, R, Bash, SQL, Snakemake
        \end{onecolentry}

        \vspace{0.2 cm}

        \begin{onecolentry}
            \textbf{Languages:} English (proficient), Chinese (native), Danish (Basic)
        \end{onecolentry}

        \vspace{0.2 cm}

        \begin{onecolentry}
            \textbf{Soft Skills:} Communication, Teamwork, Client Engagement, Multitasking, Critical Thinking, Strong Work Ethic
        \end{onecolentry}

        \vspace{0.2 cm}

        \begin{onecolentry}
            \textbf{Hobbies:} 
            I’ve been playing \textbf{badminton} since an early age—it’s a hobby that keeps me focused and energized. During holidays, I enjoy \textbf{traveling} around Europe to experience new food and cultures. I also love \textbf{hiking}, both for the chance to connect with nature and for the challenge of seeing hidden scenery that isn’t easily accessible. After moving to Denmark, I discovered a new passion: \textbf{sailing}. Last summer, we sailed to Ven and Hallands Väderö, and it was a completely new experience, and I felt free to ride the infinite waves.
        \end{onecolentry}


    
\end{document}
